\section{Introduzione al Caso di Studio}

Per illustrare concretamente come un sistema IoT possa rispondere ai requisiti e ai vincoli progettuali individuati nei capitoli precedenti, si è scelto di analizzare il progetto \textbf{SWAMP} (\textit{Smart Water Management Platform}), una piattaforma IoT per la gestione intelligente dell'irrigazione in agricoltura di precisione, sviluppata nell'ambito di una collaborazione internazionale tra Unione Europea e Brasile, finanziata dal programma Horizon 2020 \cite{kamienski2019swamp}.

L'obiettivo principale della piattaforma SWAMP è la riduzione dello spreco idrico nella produzione agricola attraverso l'automazione e l'ottimizzazione dei processi irrigui, basandosi su dati raccolti in tempo reale da reti di sensori distribuiti sul campo. Il progetto è stato validato mediante siti pilota situati sia in contesti climatici mediterranei (Italia e Spagna) sia tropicali (Brasile), rendendolo un caso di studio particolarmente rappresentativo delle sfide operative tipiche del dominio agricolo.

La scelta di questo caso è motivata dalla completezza dell'architettura proposta: SWAMP copre l'intero ciclo, dalla raccolta dei dati ambientali tramite sensori, alla trasmissione a basso consumo energetico, fino all'elaborazione in cloud e alla generazione di raccomandazioni irrigue automatizzate. Tale ampiezza consente di valutare il sistema alla luce di ciascuno dei requisiti ingegneristici precedentemente definiti.


\section{Architettura del Sistema}

La piattaforma SWAMP adotta un’architettura a più livelli, coerente con il paradigma IoT tipico delle applicazioni agricole. Sebbene l'organizzazione del sistema sia concettualmente riconducibile ai classici strati funzionali dell'IoT (percezione, rete, elaborazione intermedia e cloud), il progetto struttura formalmente l'architettura su cinque livelli specifici:

\begin{itemize}
	\item \textbf{Layer 1: Device \& Communication (Dispositivi e Comunicazione)}: Questo livello astrae una varietà di tecnologie di sensori e attuatori impiegati per acquisire parametri relativi al suolo (es. umidità), alle colture (es. stadio di crescita) e alle condizioni meteorologiche (es. temperatura dell'aria). Include l'utilizzo di sonde multiparametriche, droni per l'acquisizione di immagini e l'impiego di tecnologie di comunicazione LPWAN come LoRaWAN.
	

	\item \textbf{Layer 2: Data Acquisition, Security \& Management (Acquisizione Dati, Sicurezza e Gestione)}: Caratterizzato da protocolli e componenti software per l'acquisizione dei dati, come MQTT e LoRa Server. Questo livello include funzioni di sicurezza, gestione dei dispositivi e il componente FIWARE IoT Agent, che si occupa di tradurre la rappresentazione interna dei dati da e verso i dispositivi sul campo.
	
	\item \textbf{Layer 3: Data Management (Gestione dei Dati)}: Contiene i componenti software incaricati dell'archiviazione, elaborazione e distribuzione dei dati, basati sugli "Enablers" di FIWARE (come Orion, Quantum Leap, Cygnus e Cosmos) e sul motore semantico SEPA SPARQL. Questo strato si appoggia su un'infrastruttura distribuita composta da server cloud e nodi fog che lavorano in sinergia per gestire enormi quantità di dati.
	
	\item \textbf{Layer 4: Water Irrigation \& Distribution Models (Modelli di Irrigazione e Distribuzione Idrica)}: A questo livello risiedono i modelli agricoli tradizionali per la stima del fabbisogno idrico, basati sia sull'analisi del suolo sia sull'acquisizione di immagini tramite droni. Include inoltre modelli di ottimizzazione per la distribuzione dell'acqua e l'impiego di tecniche di intelligenza computazionale, come il machine learning.
	
	\item \textbf{Layer 5: Water Application Services (Servizi Applicativi per l'Acqua)}: Costituisce il livello più alto e fornisce i servizi finali che traducono l'elaborazione dei dati in azioni concrete, rendendole accessibili agli agricoltori e ai distributori d'acqua tramite apposite interfacce utente.
	
	
\end{itemize}


\section{Risposta ai Requisiti di Connettività e Latenza}

La piattaforma SWAMP affronta il problema della connettività in ambienti rurali attraverso l'adozione del protocollo LoRaWAN, una tecnologia LPWAN (\textit{Low Power Wide Area Network}) progettata specificamente per scenari in cui è necessaria un'ampia copertura geografica con volumi di dati contenuti.

I nodi sensore trasmettono pacchetti di dimensioni ridotte, contenenti misurazioni ambientali campionate a intervalli regolari (tipicamente ogni 10 minuti). Tale frequenza è coerente con la dinamica dei processi agronomici monitorati, che non richiedono aggiornamenti in tempo reale stretto. Il sistema opera pertanto con un modello di comunicazione \textbf{asincrono}, dove i dati vengono raccolti periodicamente e inoltrati ai gateway senza vincoli di latenza deterministica.

La copertura radio dei gateway LoRaWAN consente di servire nodi distribuiti su diversi ettari con un numero limitato di punti di accesso, riducendo la complessità infrastrutturale. Nei siti pilota, un singolo gateway è stato in grado di coprire l'intera area sperimentale, confermando l'adeguatezza della tecnologia per il contesto agricolo.



\section{Gestione dei Vincoli Energetici}

Il vincolo energetico è una criticità fondamentale affrontata nella piattaforma SWAMP, sia per quanto riguarda l'infrastruttura di rete sia per i dispositivi distribuiti sul campo. % [cite: 238, 343]

Per i nodi di \textit{sensing}, il profilo energetico del sistema è ottimizzato tramite l'adozione di tecnologie di comunicazione LPWAN, come LoRaWAN. % [cite: 188] 
Per estendere ulteriormente l'autonomia operativa, specialmente in contesti di agricoltura intensiva come il sito pilota di MATOPIBA in Brasile, il progetto prevede la sperimentazione di soluzioni per sensori wireless a bassissimo consumo energetico (\textit{ultra-low-power}) basate sull'integrazione di \textit{wake-up radio} combinate con moduli LoRaWAN. % [cite: 355] 
Questo approccio permette ai nodi di ottimizzare le fasi di trasmissione, abbattendo drasticamente i consumi. % [cite: 355]

Inoltre, l'architettura SWAMP gestisce i vincoli energetici a livello di topologia di rete attraverso l'impiego dei \textit{Field Fog Node} (FFN). % [cite: 342, 343] 
Oltre a disporre di una fonte di alimentazione stabile, questi nodi periferici agiscono come punti di aggregazione per i dispositivi LPWAN. % [cite: 343] 
Centralizzando la raccolta dei dati a livello locale prima dell'inoltro verso il cloud o verso il \textit{Fog Hub}, si ottimizza il carico trasmissivo sui singoli sensori, contribuendo all'efficienza complessiva del sistema. % [cite: 343, 345]


\section{Topologia e Ambiente Operativo}

Dal punto di vista topologico, la piattaforma SWAMP sfrutta le reti LoRaWAN, dove i nodi sensore comunicano con i gateway centrali. % [cite: 188, 345]
Questi ultimi, a loro volta, inoltrano i dati verso i livelli \textit{fog} e \textit{cloud} tramite connessioni IP, come Wi-Fi o reti cellulari (4G/5G). % [cite: 446, 448]

A differenza di quanto si possa ipotizzare per una rete agricola, i dispositivi \textbf{non sono esclusivamente statici}. Sebbene i sensori per il suolo siano installati in posizioni fisse (organizzate in zone di gestione basate sulle specifiche proprietà del terreno), il sistema integra anche nodi mobili come i droni, utilizzati per l'acquisizione di immagini per i modelli di irrigazione. % [cite: 190, 350, 354]
Inoltre, gli stessi gateway (\textit{Field Fog Node}) possono essere ospitati su infrastrutture in movimento, come i grandi sistemi di irrigazione a pivot centrale. % [cite: 344]

L'ambiente operativo dei siti pilota presenta una notevole variabilità. Il sistema è stato progettato per adattarsi a diverse configurazioni, paesi, climi, suoli e colture. % [cite: 52]
I test includono infatti scenari radicalmente diversi tra loro: dal clima della savana (\textit{cerrado}) del sito di MATOPIBA in Brasile, ai vigneti di altitudine nella regione della Mantiqueira, fino alle aree semi-aride con problemi di scarsità d'acqua a Cartagena in Spagna. % [cite: 210, 241, 267]

La validazione attraverso i quattro siti pilota distribuiti tra Brasile, Italia e Spagna ha permesso di verificare la flessibilità e la robustezza della soluzione in queste diverse condizioni ambientali reali, confermando la praticabilità dell'approccio nelle diverse realtà agricole. % [cite: 51, 52]


\section{Sicurezza e Protezione dei Dati}

Sul piano della sicurezza, gli autori del progetto SWAMP riconoscono che i sistemi IoT per l'agricoltura presentano requisiti fondamentali e complessi, tra cui la garanzia di privacy, confidenzialità, integrità, autenticazione, autorizzazione e accounting. Lo sviluppo di piattaforme IoT di successo deve necessariamente tenere conto delle sfide significative poste dalle potenziali minacce informatiche.

All'interno dell'architettura a livelli di SWAMP, le funzioni di sicurezza e di gestione dei dispositivi sono collocate nel \textit{Layer 2} (Data Acquisition, Security \& Management). Questo livello, che si interpone tra i dispositivi fisici e la gestione dei dati in cloud/fog, integra componenti specifici dedicati a questo scopo, identificati nell'architettura come moduli \textit{FIWARE Security}.

Tuttavia, il documento di riferimento specifica esplicitamente che un'analisi dettagliata dei meccanismi di protezione, dei protocolli crittografici e delle vulnerabilità fisiche esula dagli scopi specifici di quella trattazione. Di conseguenza, la piattaforma prevede un'infrastruttura architetturale predisposta per la sicurezza, delegando però i dettagli implementativi (come i meccanismi di crittografia end-to-end o la gestione avanzata delle chiavi di sessione) all'integrazione con gli standard di settore.


\section{Valutazione Critica del Caso di Studio}

Il progetto SWAMP rappresenta un esempio significativo di come un'architettura IoT completa possa essere progettata per rispondere alle reali esigenze dell'agricoltura di precisione. % [cite: 30] 
I test condotti nei siti pilota hanno permesso di validare la fattibilità tecnica della soluzione, concentrandosi in particolare sull'analisi delle prestazioni e sui limiti architetturali dell'infrastruttura software. % [cite: 31, 62]

Tuttavia, è opportuno evidenziare alcuni aspetti critici emersi dall'analisi:

\begin{itemize}
    \item \textbf{Scalabilità software}: l'incremento del numero di nodi (testato fino a 45.000 sensori simulati) ha evidenziato colli di bottiglia critici nella piattaforma FIWARE. % [cite: 638, 639, 659, 746] 
    In particolare, l'elaborazione di un elevato volume di messaggi causa un uso intensivo della CPU da parte del database MongoDB e problemi di saturazione della memoria (RAM) nel componente IoT Agent, portando a instabilità del sistema. % [cite: 64, 65, 691, 695, 727, 728]
    
    \item \textbf{Dipendenza dalla connettività cloud}: sebbene il livello \textit{fog} garantisca una certa autonomia locale e robustezza contro le disconnessioni, le funzionalità più avanzate (modelli agronomici complessi, \textit{machine learning}) dipendono dalla disponibilità della connessione Internet, che in contesti rurali (come il sito pilota brasiliano) può risultare intermittente o instabile. % [cite: 346, 348, 350, 403, 405]
    
    \item \textbf{Complessità di Deployment}: non esiste una configurazione universale per l'IoT agricolo. % [cite: 737] 
    L'adattamento della piattaforma ai diversi requisiti di connettività, infrastruttura e obiettivi (irrigazione o distribuzione idrica) richiede la progettazione di componenti altamente specifici e personalizzati per la singola azienda agricola, incrementando lo sforzo di sviluppo. % [cite: 52, 56, 57, 204, 207]
    
    \item \textbf{Gestione della semantica dei dati}: l'adozione del framework FIWARE favorisce la standardizzazione, ma per gestire scenari più complessi e interconnessi (come l'ottimizzazione della distribuzione idrica consortile) si è resa necessaria l'integrazione di motori semantici aggiuntivi (come SEPA SPARQL), aumentando la complessità dell'architettura per tradurre i modelli di dati. % [cite: 488, 489, 492, 749, 750, 752]
\end{itemize}

\noindent Nel complesso, il caso SWAMP dimostra che i requisiti individuati nell'analisi teorica trovano riscontro concreto nelle sfide di un sistema reale. % [cite: 731] 
Le soluzioni adottate per la gestione della connettività rurale, l'integrazione \textit{cloud-fog} e la standardizzazione dei dati risultano coerenti con i vincoli tipici del dominio agricolo, evidenziando al contempo come la maturità delle piattaforme IoT richieda ancora ottimizzazioni per gestire elaborazioni massive su larga scala. % [cite: 59, 745, 746, 758]